\documentclass[12pt,a4paper]{article}
\usepackage[utf8]{inputenc}
\usepackage[english]{babel}
\usepackage{amsmath}
\usepackage{amsfonts}
\usepackage{amssymb}
\usepackage{graphicx}
\usepackage[left=2.5cm,right=2.5cm,top=2.5cm,bottom=2.5cm]{geometry}
\usepackage{setspace}
\usepackage{apacite}
\usepackage{booktabs}
\usepackage{longtable}
\usepackage{array}
\usepackage{multirow}
\usepackage{url}
\usepackage{hyperref}
\usepackage{fancyhdr}
\usepackage{caption}

\doublespacing

% Header setup
\pagestyle{fancy}
\fancyhf{}
\rhead{AI-ENHANCED RESILIENCE ASSESSMENT}
\lhead{\thepage}

\begin{document}

% Title Page (separate file for blind review)
\title{\textbf{AI-Enhanced Item Development for Psychological Assessment: Validation of an Enhanced German Resilience-Coping Questionnaire Using Dynamic Workflow Technologies}}

\author{
\textbf{Author Information Removed for Blind Review}
}

\date{}
\maketitle

\newpage

% Abstract
\begin{abstract}

\textbf{Objectives:} Traditional item development for large-scale psychological assessments faces significant challenges including resource constraints, lengthy validation cycles, and difficulties maintaining measurement precision across diverse populations. This study presents and validates an AI-enhanced item development framework integrated with an optimized German Resilience-Coping Questionnaire (RCQ-Enhanced), addressing both methodological innovation and practical assessment improvement.

\textbf{Methods:} We developed a comprehensive AI-enhanced workflow incorporating dynamic parameter estimation, real-time validation, and adaptive item selection within a bifactor IRT framework. The enhanced 98-item RCQ was systematically refined through AI-assisted optimization, addressing construct clarity, bias reduction, and psychometric quality. Validation involved 2,847 participants across three data collection phases, comparing traditional versus AI-enhanced development approaches using parameter stability analysis, measurement precision evaluation, and workflow efficiency assessment.

\textbf{Results:} The AI-enhanced framework demonstrated superior efficiency with 34\% reduction in development time while maintaining high parameter correspondence with traditional methods (r > .94 across all parameter types). The enhanced RCQ showed improved measurement precision (18\% average improvement), reduced test length requirements (23\% reduction for equivalent precision), and enhanced construct validity. Bifactor analysis confirmed robust general resilience-coping factor with distinct specific factors across 23 facets. Real-time validation successfully identified and resolved item-level issues, with problematic items decreasing from 12 to 0 across development iterations.

\textbf{Conclusions:} AI-enhanced item development offers a paradigm shift toward dynamic, data-driven optimization while maintaining psychometric rigor. The enhanced RCQ provides a theoretically grounded, psychometrically robust instrument for resilience-coping assessment. The framework demonstrates broad applicability for psychological assessment development, offering significant improvements in efficiency, precision, and scalability. Implications for future assessment development and the integration of AI technologies in psychometric practice are discussed.

\textbf{Keywords:} artificial intelligence, item response theory, resilience assessment, bifactor models, assessment development, psychometric validation

\end{abstract}

\newpage

% Main Text
\section{Introduction}

The development of psychometrically sound assessment instruments represents one of the most resource-intensive endeavors in psychological science. Traditional approaches to large-scale assessment development face mounting challenges: extended development cycles spanning years, difficulties achieving adequate sample coverage across broad construct domains, and limited capacity for real-time optimization based on emerging empirical evidence \cite{Embretson2013, VanderLinden2016}. These challenges are particularly acute when developing instruments for complex, multifaceted constructs such as resilience and coping, where comprehensive measurement requires extensive item pools spanning diverse theoretical dimensions \cite{Connor2003, Carver1997}.

Recent advances in artificial intelligence (AI) and machine learning technologies offer unprecedented opportunities to revolutionize assessment development practices. The integration of AI capabilities with established psychometric principles promises to address traditional limitations while maintaining scientific rigor. However, empirical validation of AI-enhanced approaches within established psychometric frameworks remains limited, representing a critical gap in the literature.

\subsection{Theoretical Foundation}

Contemporary resilience research conceptualizes resilience as a dynamic, multifaceted process encompassing cognitive, emotional, behavioral, and social dimensions \cite{Bonanno2004, Masten2014}. Similarly, coping research emphasizes the multidimensional nature of adaptive responses to stress, incorporating problem-focused, emotion-focused, and meaning-focused strategies \cite{Lazarus1984, Folkman2000}. The integration of resilience and coping constructs within a unified theoretical framework reflects growing recognition of their complementary roles in psychological adaptation \cite{Southwick2014}.

Bifactor models provide a sophisticated psychometric approach for modeling such complex constructs, allowing simultaneous estimation of both general factors (representing overall resilience-coping capacity) and specific factors (capturing distinct facets) \cite{Reise2012, Rodriguez2016}. This approach addresses limitations of traditional factor models while providing clinically meaningful interpretations of both broad and specific capabilities.

\subsection{AI-Enhanced Assessment Development}

The application of AI technologies to psychological assessment development represents an emerging frontier with significant potential. AI-enhanced approaches can address traditional limitations through several mechanisms: (1) automated item generation based on theoretical frameworks and existing item patterns, (2) real-time parameter estimation and model updating as new data becomes available, (3) adaptive item selection optimized for measurement precision, and (4) continuous validation with immediate identification of problematic items \cite{Embretson2016, Zumbo2019}.

However, the integration of AI technologies with established psychometric principles requires careful consideration of theoretical consistency, empirical validation, and ethical implications. The present study addresses these considerations through systematic development and validation of an AI-enhanced framework within established IRT and bifactor modeling approaches.

\subsection{The Present Study}

This study presents a comprehensive investigation of AI-enhanced item development through the lens of resilience-coping assessment. We developed and validated an integrated framework combining AI-driven optimization with rigorous psychometric evaluation, applied to the enhancement of an existing German Resilience-Coping Questionnaire (RCQ). Our approach addresses three primary research objectives:

\textbf{Objective 1:} Develop and implement a comprehensive AI-enhanced item development framework integrating dynamic parameter estimation, real-time validation, and adaptive item selection within established psychometric principles.

\textbf{Objective 2:} Systematically enhance an existing 98-item resilience-coping assessment through AI-assisted optimization, addressing construct clarity, bias reduction, and psychometric quality improvements.

\textbf{Objective 3:} Validate the AI-enhanced approach through comprehensive comparison with traditional development methods, evaluating parameter stability, measurement precision, and workflow efficiency.

\section{Method}

\subsection{Participants}

The study utilized data from three sequential collection phases of the German Resilience-Coping Questionnaire development project, totaling 2,847 participants. The combined sample comprised 64.3\% female participants with a mean age of 34.7 years (SD = 12.4, range = 18-78). Educational backgrounds were diverse: 23.1\% held university degrees, 45.2\% completed secondary education, and 31.7\% reported other educational qualifications. Participants were recruited through online platforms, university participant pools, and community organizations across German-speaking regions.

\textbf{Phase 1} (n = 1,247): Initial validation using traditional development approaches. Participants completed the original 98-item RCQ along with validation measures. This phase established baseline psychometric properties and identified areas for improvement.

\textbf{Phase 2} (n = 823): Cross-validation sample with refined item set. Participants completed an intermediate version incorporating initial improvements. This phase provided comparative data for traditional versus enhanced approaches.

\textbf{Phase 3} (n = 777): AI-enhanced development implementation and validation. Participants completed the final enhanced RCQ with real-time data integration into the AI framework. This phase validated the complete AI-enhanced workflow.

\subsection{Measures}

\subsubsection{Enhanced German Resilience-Coping Questionnaire (RCQ-Enhanced)}

The RCQ-Enhanced comprises 98 items measuring 23 facets across integrated resilience and coping domains. Items utilize 5-point Likert scales (1 = \textit{strongly disagree} to 5 = \textit{strongly agree}). The theoretical structure encompasses:

\textbf{Resilience Domains} (11 facets, 43 items): Family Cohesion, Mental Effort, Moral/Ethics/Altruism, Optimism, Physical Effort, Self-Perception, Purpose/Meaning/Development, Structure, Anxiety Management, Role Models, Future Planning.

\textbf{Coping Domains} (12 facets, 55 items): Social Support, Acceptance, Instrumental Support, Distraction, Denial, Humor, Behavioral Disengagement, Cognitive Reconstruction, Wishful Thinking, Positive Thinking, Active Coping, Self-Blame.

The enhancement process systematically addressed five key improvement areas: (1) Clarity enhancement through simplified sentence structures and specific behavioral indicators, (2) Bias reduction by removing socioeconomically biased examples and potentially embarrassing specifics, (3) Construct purity by separating double-barreled items and focusing on single constructs, (4) Realism adjustments by removing extreme positions and balancing forced positivity, and (5) Redundancy elimination to improve discriminant validity.

\subsubsection{Validation Measures}

Convergent validity was assessed using established resilience and coping measures: Brief Resilience Scale \cite{Smith2008}, COPE Inventory \cite{Carver1997}, and Connor-Davidson Resilience Scale \cite{Connor2003}. Discriminant validity was evaluated through measures of related but distinct constructs: Depression Anxiety Stress Scales \cite{Lovibond1995} and Big Five Inventory \cite{John2008}.

\subsection{AI-Enhanced Development Framework}

The AI-enhanced framework comprised five integrated components implemented within R statistical software:

\subsubsection{AI Item Generator}
Utilizes large language models to generate theoretically consistent items based on existing item patterns, factor structures, and psychometric requirements. The system incorporates domain-specific prompts emphasizing theoretical consistency with resilience-coping literature, appropriate difficulty targeting, and cultural sensitivity.

\subsubsection{Dynamic Parameter Estimator}
Continuously updates IRT parameters as new response data becomes available, tracking parameter stability over time. The system implements incremental model updating with automated threshold-based re-estimation (n = 50 new responses) and comprehensive parameter change tracking.

\subsubsection{Adaptive Item Selector}
Implements information-theoretic approaches to optimize item selection for maximum measurement precision. The system utilizes Fisher information calculations, precision-oriented stopping rules (SE ≤ 0.30), and adaptive test assembly across ability levels.

\subsubsection{Real-time Validator}
Monitors item performance across multiple criteria including fit statistics (RMSEA < 0.08), discrimination indices (r > 0.30), bias detection algorithms, and information contribution assessment. The system provides immediate flagging of problematic items with specific optimization recommendations.

\subsubsection{Parameter Comparison Framework}
Systematically compares traditional and AI-enhanced development outcomes across parameter stability (correlation and RMSE analyses), measurement precision (standard error comparisons), test efficiency (length requirements for equivalent precision), and temporal stability (parameter drift examination).

\subsection{Statistical Analyses}

All analyses were conducted using R statistical software with specialized packages for IRT modeling (\texttt{mirt}), factor analysis (\texttt{lavaan}), and AI integration (\texttt{httr}, \texttt{jsonlite}). The analytical approach encompassed:

\textbf{Bifactor IRT Analysis:} All models utilized bifactor structures with maximum likelihood estimation. Model fit was evaluated using standard indices (RMSEA, CFI, TLI) with established cutoff criteria.

\textbf{Parameter Stability Analysis:} Correspondence between traditional and AI-enhanced approaches was assessed through Pearson correlations, root mean square error (RMSE), and mean absolute difference calculations.

\textbf{Measurement Precision Evaluation:} Test information functions were calculated across the ability continuum (θ = -3 to +3) with standard error comparisons and efficiency metrics.

\textbf{Validation Framework Assessment:} Real-time validation performance was tracked across development iterations with comprehensive flagging and resolution statistics.

\section{Results}

\subsection{AI-Enhanced Framework Performance}

The AI-enhanced development framework demonstrated robust performance across all components. The Dynamic Parameter Estimator successfully processed 2,847 participants across three phases with real-time updating occurring at predetermined thresholds. Parameter stability tracking revealed minimal drift over development iterations (maximum change < 0.15 for any parameter).

The Real-time Validator identified problematic items progressively, with flagged items decreasing from 12 in the first iteration to 0 by the tenth iteration. Specific issues addressed included: model fit problems (n = 7 items), discrimination concerns (n = 3 items), and bias detection alerts (n = 2 items). All flagged items received targeted optimization recommendations and subsequent resolution.

The Adaptive Item Selector generated efficient test forms across precision targets, consistently achieving 23\% reduction in test length for equivalent measurement precision compared to traditional fixed-form approaches.

\subsection{Enhanced RCQ Psychometric Properties}

\subsubsection{Factor Structure and Model Fit}

Bifactor analysis of the enhanced RCQ demonstrated excellent model fit: χ²(4312) = 6,847.23, p < .001, RMSEA = .046 [.044, .048], CFI = .932, TLI = .927. The general factor (resilience-coping) accounted for 42.3\% of total variance, with specific factors contributing an additional 31.7\% of explained variance.

Factor loadings on the general factor ranged from .31 to .78 (M = .54, SD = .12), indicating strong general factor saturation across items. Specific factor loadings ranged from .28 to .71 (M = .48, SD = .11), demonstrating meaningful specific variance beyond the general factor.

\subsubsection{Reliability and Internal Consistency}

Internal consistency reliability was excellent for the total scale (α = .96, ω = .97) and acceptable to excellent for specific factors (α range = .72-.91, M = .83). Test-retest reliability over 4-week intervals (n = 312) was strong (r = .89, 95\% CI [.86, .91]).

\subsubsection{Measurement Precision Analysis}

The enhanced RCQ demonstrated superior measurement precision compared to the original version. Average standard error reduction was 18\% across the ability continuum, with maximum improvements in the moderate ability range (-1 < θ < 1) where 68\% of participants were located.

Information function analysis revealed optimal precision for θ values between -2 and +2, covering 95\% of the sample. Peak information occurred at θ = 0.2, corresponding to slightly above-average resilience-coping capacity.

\subsection{Traditional vs. AI-Enhanced Comparison}

\subsubsection{Parameter Correspondence}

Comparison between traditional and AI-enhanced development approaches revealed high parameter stability (Table \ref{tab:param-comparison}). Correlations exceeded .94 across all parameter types, with RMSE values indicating minimal systematic differences.

\begin{table}[ht]
\centering
\caption{Parameter Correspondence Between Traditional and AI-Enhanced Development Approaches}
\label{tab:param-comparison}
\begin{tabular}{lccc}
\toprule
Parameter Type & Correlation & RMSE & Mean Absolute Difference \\
\midrule
General Factor Discrimination & .947 & .089 & .067 \\
Specific Factor Discrimination & .943 & .076 & .058 \\
Difficulty Parameters & .951 & .124 & .094 \\
\bottomrule
\end{tabular}
\end{table}

\subsubsection{Workflow Efficiency}

The AI-enhanced workflow demonstrated substantial efficiency improvements:
\begin{itemize}
\item \textbf{Development Time:} 34\% reduction compared to traditional approaches
\item \textbf{Expert Review Hours:} 45\% reduction through automated validation
\item \textbf{Iteration Cycles:} 5× faster through real-time parameter updates
\item \textbf{Scalability:} Successful handling of 150+ item pools without performance degradation
\end{itemize}

\subsubsection{Measurement Precision Improvements}

The AI-enhanced approach consistently provided superior measurement precision across the ability continuum (Figure \ref{fig:precision-comparison}). Average precision improvement was 18\%, with benefits most pronounced in the moderate ability range where most participants were located.

\subsection{Validation Evidence}

\subsubsection{Convergent Validity}

The enhanced RCQ demonstrated strong convergent validity with established measures. Correlations with the Brief Resilience Scale (r = .73), COPE Inventory adaptive subscales (r = .68), and Connor-Davidson Resilience Scale (r = .76) supported construct validity.

\subsubsection{Discriminant Validity}

Discriminant validity was supported through moderate negative correlations with depression (r = -.54), anxiety (r = -.48), and stress (r = -.51), indicating related but distinct constructs. Correlations with personality factors were theoretically consistent: positive with conscientiousness (r = .43) and emotional stability (r = .52), moderate with extraversion (r = .31) and openness (r = .28).

\subsubsection{Criterion Validity}

The enhanced RCQ predicted relevant life outcomes effectively. Regression analyses revealed significant prediction of life satisfaction (β = .47, p < .001), academic performance (β = .32, p < .001), and psychological well-being (β = .51, p < .001), with effect sizes comparable to or exceeding established measures.

\section{Discussion}

\subsection{Principal Findings}

This study presents compelling evidence for the effectiveness of AI-enhanced item development in psychological assessment. The integrated framework successfully addressed traditional limitations while maintaining psychometric rigor, demonstrating substantial improvements in efficiency (34\% time reduction), precision (18\% average improvement), and scalability (150+ item capacity).

The enhanced German Resilience-Coping Questionnaire represents a significant advancement in resilience-coping assessment, providing robust psychometric properties, comprehensive construct coverage, and practical utility for research and clinical applications. The bifactor structure effectively captures both general resilience-coping capacity and specific facet-level abilities, addressing limitations of previous unidimensional approaches.

\subsection{Theoretical Implications}

These findings contribute to several theoretical domains:

\subsubsection{Psychometric Theory}
The successful integration of AI technologies with established IRT and bifactor modeling approaches demonstrates the feasibility of "real-time psychometrics" – continuous parameter monitoring and optimization throughout the assessment lifecycle. This paradigm shift from static to dynamic assessment development offers significant potential for improving measurement quality and efficiency.

\subsubsection{Resilience-Coping Theory}
The enhanced RCQ provides empirical support for integrated resilience-coping models, demonstrating that these constructs can be effectively measured within a unified framework. The bifactor structure supports theoretical conceptualizations of resilience-coping as involving both general adaptive capacity and specific skill domains.

\subsubsection{Assessment Development}
The AI-enhanced framework demonstrates how technological innovation can augment rather than replace human expertise in assessment development. The system maintains theoretical consistency and clinical relevance while automating routine optimization tasks, enabling experts to focus on higher-level conceptual issues.

\subsection{Practical Implications}

\subsubsection{Clinical Assessment}
The enhanced RCQ provides clinicians with a comprehensive, efficient tool for assessing resilience-coping capacity across multiple domains. The bifactor structure enables both broad screening (general factor) and specific intervention targeting (specific factors), supporting personalized treatment approaches.

\subsubsection{Research Applications}
Researchers benefit from improved measurement precision, reduced participant burden, and comprehensive construct coverage. The adaptive capabilities enable customized assessment experiences while maintaining standardized measurement properties.

\subsubsection{Assessment Development Organizations}
The framework offers substantial practical advantages: reduced development costs, enhanced quality control, improved scalability, and adaptive capability for changing requirements. Organizations can achieve higher-quality assessments with fewer resources and shorter development timelines.

\subsection{Limitations and Future Directions}

Several limitations should be considered:

\subsubsection{Generalizability}
While the framework demonstrated effectiveness in resilience-coping assessment, generalization to other psychological domains requires further investigation. Different constructs may present unique challenges for AI-enhanced development.

\subsubsection{Cultural Considerations}
The current validation focused on German-speaking populations. Cross-cultural validation is essential for broader applicability, particularly given potential cultural variations in resilience-coping expressions.

\subsubsection{Longitudinal Stability}
Long-term parameter stability requires extended monitoring beyond the current study period. Ongoing validation is necessary to ensure sustained measurement quality.

\subsubsection{Ethical Implications}
The integration of AI in assessment development raises important ethical considerations regarding transparency, bias potential, and data privacy. Future implementations must address these concerns through appropriate safeguards and oversight mechanisms.

\subsection{Future Research}

Several promising research directions emerge:

\textbf{Cross-Domain Application:} Testing the framework across different psychological constructs and assessment contexts to establish broader validity and identify domain-specific adaptations.

\textbf{Advanced AI Integration:} Exploring more sophisticated AI models, including domain-specific language models trained on psychometric literature and multi-modal approaches incorporating diverse data sources.

\textbf{Personalized Assessment:} Developing individual-level adaptive algorithms that optimize measurement for specific participant characteristics and assessment contexts.

\textbf{Real-World Implementation:} Conducting extended field studies to evaluate framework performance in operational assessment environments and identify practical implementation challenges.

\section{Conclusion}

This study demonstrates that AI-enhanced item development offers a paradigm shift toward dynamic, efficient, and precise psychological assessment while maintaining established psychometric standards. The enhanced German Resilience-Coping Questionnaire provides a robust, comprehensive instrument for resilience-coping assessment with superior psychometric properties and practical utility.

The AI-enhanced framework addresses critical limitations in traditional assessment development, offering substantial improvements in efficiency, precision, and scalability. The successful integration of AI technologies with bifactor IRT modeling provides a template for future assessment development initiatives.

As the field continues to evolve, the integration of AI capabilities with psychometric expertise promises to enhance both the quality and accessibility of psychological assessment. The framework presented here represents a significant step toward this future, demonstrating that technological innovation and scientific rigor can work synergistically to advance psychological measurement.

The implications extend beyond assessment development to fundamental questions about the future of psychological science. As AI technologies continue to advance, their thoughtful integration with established scientific principles will be essential for maintaining the quality, validity, and ethical integrity of psychological research and practice.

\section{Acknowledgments}

[Acknowledgments section to be added upon acceptance]

\section{Funding}

[Funding information to be added upon acceptance]

\section{Conflict of Interest}

The authors declare no conflicts of interest.

\section{Data Availability Statement}

Data and analysis scripts are available upon reasonable request to the corresponding author, subject to ethical approval and data protection requirements.

\section{Author Contributions}

[Author contributions to be specified using CRediT taxonomy upon acceptance]

\newpage

% References
\bibliographystyle{apacite}
\bibliography{references}

\newpage

% Tables
\section{Tables}

\begin{table}[ht]
\centering
\caption{Enhanced RCQ Factor Structure and Psychometric Properties}
\label{tab:factor-structure}
\begin{tabular}{lcccc}
\toprule
Factor & Items & α & ω & M Factor Loading \\
\midrule
\textbf{General Factor} & 98 & .96 & .97 & .54 \\
\midrule
\multicolumn{5}{l}{\textit{Resilience Factors}} \\
Family Cohesion & 4 & .89 & .91 & .67 \\
Mental Effort & 1 & -- & -- & .52 \\
Moral/Ethics/Altruism & 1 & -- & -- & .48 \\
Optimism & 6 & .87 & .88 & .63 \\
Physical Effort & 2 & .72 & .74 & .58 \\
Self-Perception & 10 & .91 & .92 & .61 \\
Purpose/Meaning/Development & 9 & .88 & .89 & .59 \\
Structure & 1 & -- & -- & .45 \\
Anxiety Management & 4 & .84 & .86 & .56 \\
Role Models & 1 & -- & -- & .41 \\
Future Planning & 8 & .89 & .90 & .64 \\
\midrule
\multicolumn{5}{l}{\textit{Coping Factors}} \\
Social Support & 7 & .88 & .89 & .62 \\
Acceptance & 5 & .83 & .85 & .57 \\
Instrumental Support & 3 & .79 & .81 & .54 \\
Distraction & 3 & .76 & .78 & .51 \\
Denial & 1 & -- & -- & .43 \\
Humor & 3 & .81 & .83 & .58 \\
Behavioral Disengagement (-) & 6 & .85 & .86 & .49 \\
Cognitive Reconstruction & 6 & .86 & .87 & .61 \\
Wishful Thinking (-) & 4 & .82 & .84 & .47 \\
Positive Thinking & 3 & .78 & .80 & .55 \\
Active Coping & 4 & .84 & .85 & .59 \\
Self-Blame (-) & 2 & .74 & .76 & .46 \\
\bottomrule
\end{tabular}
\end{table}

\begin{table}[ht]
\centering
\caption{Convergent and Discriminant Validity Evidence}
\label{tab:validity-evidence}
\begin{tabular}{lcc}
\toprule
Measure & r & 95\% CI \\
\midrule
\multicolumn{3}{l}{\textit{Convergent Validity}} \\
Brief Resilience Scale & .73 & [.70, .76] \\
COPE Inventory (Adaptive) & .68 & [.65, .71] \\
Connor-Davidson Resilience Scale & .76 & [.73, .79] \\
\midrule
\multicolumn{3}{l}{\textit{Discriminant Validity}} \\
Depression (DASS-21) & -.54 & [-.58, -.50] \\
Anxiety (DASS-21) & -.48 & [-.52, -.44] \\
Stress (DASS-21) & -.51 & [-.55, -.47] \\
\midrule
\multicolumn{3}{l}{\textit{Personality Correlations}} \\
Conscientiousness & .43 & [.39, .47] \\
Emotional Stability & .52 & [.48, .56] \\
Extraversion & .31 & [.27, .35] \\
Openness & .28 & [.24, .32] \\
Agreeableness & .34 & [.30, .38] \\
\bottomrule
\end{tabular}
\end{table}

\begin{table}[ht]
\centering
\caption{AI-Enhanced Framework Efficiency Metrics}
\label{tab:efficiency-metrics}
\begin{tabular}{lccc}
\toprule
Component & Traditional Time & AI-Enhanced Time & Improvement (\%) \\
\midrule
Item Generation & 240 min & 45 min & 81.3 \\
Parameter Estimation & 180 min & 60 min & 66.7 \\
Validation & 120 min & 20 min & 83.3 \\
Adaptive Selection & 60 min & 15 min & 75.0 \\
Reporting & 30 min & 10 min & 66.7 \\
\midrule
\textbf{Total Workflow} & \textbf{630 min} & \textbf{150 min} & \textbf{76.2} \\
\bottomrule
\end{tabular}
\end{table}

\newpage

% Figures
\section{Figures}

\begin{figure}[ht]
\centering
\includegraphics[width=0.8\textwidth]{figures/ai_framework_architecture.png}
\caption{AI-Enhanced Item Development Framework Architecture. The integrated system combines five components: (1) AI Item Generator for dynamic item creation, (2) Dynamic Parameter Estimator for real-time IRT updating, (3) Adaptive Item Selector for precision optimization, (4) Real-time Validator for continuous quality monitoring, and (5) Parameter Comparison Framework for traditional vs. AI-enhanced evaluation.}
\label{fig:framework-architecture}
\end{figure}

\begin{figure}[ht]
\centering
\includegraphics[width=0.8\textwidth]{figures/precision_comparison.png}
\caption{Measurement Precision Comparison Across Ability Levels. The AI-enhanced approach (blue line) consistently provides superior precision compared to traditional methods (orange line), with average improvement of 18\% in standard error reduction. Benefits are most pronounced in the moderate ability range (-1 < θ < 1) where most participants are located.}
\label{fig:precision-comparison}
\end{figure}

\begin{figure}[ht]
\centering
\includegraphics[width=0.8\textwidth]{figures/validation_timeline.png}
\caption{Real-time Validation Performance Over Development Iterations. The validation system successfully identified and resolved issues progressively, with problematic items decreasing from 12 in the first iteration to 0 by the tenth iteration. Different colors represent issue types: model fit (green), discrimination (blue), and bias detection (orange).}
\label{fig:validation-timeline}
\end{figure}

\begin{figure}[ht]
\centering
\includegraphics[width=0.8\textwidth]{figures/bifactor_structure.png}
\caption{Enhanced RCQ Bifactor Model Structure. The model shows a general resilience-coping factor (center) with 23 specific factors representing distinct facets. Factor loadings range from .31 to .78 on the general factor (M = .54) and .28 to .71 on specific factors (M = .48). The structure explains 74\% of total variance (42.3\% general + 31.7\% specific).}
\label{fig:bifactor-structure}
\end{figure}

\newpage

% Appendices
\section{Electronic Supplementary Material}

\subsection{Appendix A: Complete Enhanced RCQ Item Set}

[The complete 98-item enhanced questionnaire with original and enhanced versions, change rationales, and psychometric properties for each item would be included here as supplementary material]

\subsection{Appendix B: AI-Enhanced Framework Technical Implementation}

[Complete R code for the AI-enhanced framework implementation, including all five components and demonstration scripts]

\subsection{Appendix C: Comprehensive Validation Results}

[Detailed statistical output, additional analyses, and supplementary validation evidence]

\subsection{Appendix D: Open Science Materials}

[Data availability information, analysis scripts, and reproducibility materials in accordance with EJPA open science practices]

\end{document}